\documentclass[11pt,oneside,openany]{book}
\usepackage[T1]{fontenc}
\usepackage[utf8]{inputenc}
\usepackage{lmodern}
\usepackage{amsmath,amssymb,amsthm,mathtools}
\usepackage{geometry,booktabs,longtable,array,enumitem,microtype,xcolor}
\usepackage[colorlinks=true,linkcolor=black,citecolor=black,urlcolor=blue!50!black,breaklinks=true]{hyperref}
\usepackage{fancyhdr}
\usepackage{titlesec}
\usepackage{tocloft}
\setlength{\cftbeforechapskip}{3pt}\setlength{\cftbeforesecskip}{0pt}
\titleclass{\chapter}{straight}
\titleformat{\chapter}[hang]{\normalfont\LARGE\bfseries\raggedright}{\chaptertitlename\ \thechapter}{0.6em}{}
\titlespacing*{\chapter}{0pt}{26pt}{14pt}
\usepackage{graphicx}
\usepackage{tikz}
\usetikzlibrary{arrows.meta,positioning,calc,decorations.pathreplacing}
\usepackage[font=small,labelfont=bf]{caption}
\usepackage{graphicx}\graphicspath{{./}{../}}
\geometry{a4paper,margin=28mm,headheight=14pt}
\setlength{\parindent}{0pt}
\setlength{\parskip}{5pt}
\setcounter{tocdepth}{2}
\pagestyle{fancy}
\fancyhf{}
\fancyhead[L]{\footnotesize Mathematical Foundations of the UNNS Substrate}
\fancyhead[R]{\footnotesize Volume II}
\fancyfoot[C]{\thepage}
\newtheorem{theorem}{Theorem}[chapter]
\newtheorem{proposition}[theorem]{Proposition}
\newtheorem{lemma}[theorem]{Lemma}
\newtheorem{corollary}[theorem]{Corollary}
\theoremstyle{definition}
\newtheorem{definition}[theorem]{Definition}
\newtheorem{example}[theorem]{Example}
\theoremstyle{remark}
\newtheorem{remark}[theorem]{Remark}
\newtheorem{problem}[theorem]{Open problem}
\newcommand{\Q}{\mathbb{Q}}
\newcommand{\N}{\mathbb{N}}
\newcommand{\Z}{\mathbb{Z}}
\newcommand{\Fx}{\mathcal{F}}
\newcommand{\Dx}{\mathcal{D}}
\newcommand{\lcmop}{\operatorname{lcm}}
\newcommand{\ord}{\operatorname{ord}}
\newcommand{\vv}{\operatorname{v}}
\newcommand{\sqf}{\operatorname{sqf}}
\newcommand{\Dcoll}{\Delta_{\mathrm{coll}}}
\newcommand{\Abs}{\operatorname{Abs}}
\newcommand{\Hcl}{H_{\mathrm{class}}}
\newcommand{\Kpoly}{K_{\mathrm{poly}}}
\newcommand{\Rbd}{\mathcal{R}}
\newcommand{\scope}[1]{\textsf{\small Scope #1.}\ }
\newcommand{\record}[1]{\textup{\textsc{Source record: #1.}}\ }
\newcommand{\rec}[1]{\record{#1}\label{th:#1}}
\newcommand{\srcpath}[1]{\href{https://github.com/ukbbi/UNNS/blob/bb8b12d753c600c597feb00df17fe139c7cc65fa/projects/UNNS_MATH_FOUNDATIONS/sources/#1}{\nolinkurl{sources/#1}}}
\newcommand{\src}[2]{\href{https://github.com/ukbbi/UNNS/blob/bb8b12d753c600c597feb00df17fe139c7cc65fa/projects/UNNS_MATH_FOUNDATIONS/sources/#1}{#2}}
\newcommand{\figpath}[1]{figs/#1}
\tikzset{box/.style={draw,rounded corners=2pt,align=center,inner sep=4pt,font=\small},
 pt/.style={circle,fill=black,inner sep=1.4pt},lab/.style={font=\footnotesize}}
\allowdisplaybreaks[1]
